\scanpage{23}
$\mathcal G$ --- lequel opère donc sur $T$. Ceci donne naissance automatiquement à
\diagram{\begin{tikzcd}[ampersand replacement=\&]
\Phi:\tau \arrow[r,dash] \& \operatorname{Hom}_{\mathrm{gr}}(T,\pi)
\end{tikzcd}}
\ldots

Ceci dit, on a un hom du groupoïde $\langle\tau,\pi\rangle$ dans le groupoïde des rev. universels de $\Pi$, constant sur les objets et de valeur $\widetilde\Pi$ sur ceux-ci, qui envoie donc l'ens des flèches dans $\pi$ --- et qui pour $t,t'\in\tau$, désignant par $\pi_{t,t'}$ \inserted{$=\operatorname{Isom}_{\langle\tau,\pi\rangle}(t,t')$} l'ens des $\lambda\in\pi$ tels que $\lambda t=t'$, est l'inclusion de $\pi_{t,t'}$ dans $\pi$. Ce foncteur
\diagram{\begin{tikzcd}[ampersand replacement=\&]
\langle\tau,\pi\rangle \arrow[r] \& \begin{gathered}\text{Cat des rev. universels}\\\text{de }\Pi\text{ égaux à }\widetilde\Pi\\\simeq(\text{un objet et groupe d'isom--- }\pi)\end{gathered}
\end{tikzcd}}
est évidemment compatible avec op. de $\mathcal E$ sur $\unsure{\Pi|\widetilde\Pi}$ donc sur $\unsure{K|\pi}\ldots$
\mspage{22}
